\begin{figure}[H]
\centering
\resizebox{\textwidth}{!}{%
\begin{tikzpicture}[
    node distance=0pt,
    adp1/.style={draw, fill=blue!25, minimum width=1.4cm, minimum height=0.6cm,
                 font=\small, align=center},
    adp2/.style={draw, fill=red!25,  minimum width=1.4cm, minimum height=0.6cm,
                 font=\small, align=center},
    primer/.style={draw, fill=orange!25, minimum width=2.0cm, minimum height=0.6cm,
                   font=\small, align=center},
    ins/.style={draw, fill=green!20, minimum width=2.8cm, minimum height=0.6cm,
                font=\small, align=center},
    lbl/.style={font=\small\itshape, anchor=east},
]

% ---- Scenario A ----
\node[lbl] at (0, 1.4) {Scenario A:};
\node[adp1] (a1)   at (1.6,  1.4) {Adp1};
\node[primer, right=of a1]  (afp)  {Fwd Primer};
\node[ins,    right=of afp] (ains) {Insert};
\node[primer, right=of ains](arp)  {Rev Primer$_{RC}$};
\node[adp1,   right=of arp] (a1rc) {Adp1$_{RC}$};
\node[font=\small, gray, anchor=west] at ([xshift=6pt]a1rc.east)
    {same adaptor on both ends $\Rightarrow$ one adaptor sequence needed};

% ---- Scenario B ----
\node[lbl] at (0, 0) {Scenario B:};
\node[adp1] (b1)   at (1.6, 0) {Adp1};
\node[primer, right=of b1]  (bfp)  {Fwd Primer};
\node[ins,    right=of bfp] (bins) {Insert};
\node[primer, right=of bins](brp)  {Rev Primer$_{RC}$};
\node[adp2,   right=of brp] (b2rc) {Adp2$_{RC}$};
\node[font=\small, anchor=west] at ([xshift=6pt]b2rc.east)
    {\textbf{different adaptors} $\Rightarrow$ two distinct sequences needed \checkmark};

\end{tikzpicture}%
}
\caption{The two possible adaptor configurations for a paired-end amplicon library.
Scenario~A uses the same adaptor on both ends (one sequence suffices);
Scenario~B uses two distinct adaptors (two sequences required).
Since two distinct adaptor sequences (Adp1 $\neq$ Adp2) were provided with this dataset,
Scenario~B was assumed and subsequently confirmed empirically
(Appendix~\ref{app:adaptor-scenarios}).}
\label{fig:adaptor-scenarios}
\end{figure}
